% ============================================================================
% CrossDisc-Bench: A Knowledge Graph-Grounded Benchmark for
% Cross-Disciplinary Scientific Hypothesis Generation
% NeurIPS 2025 Datasets & Benchmarks Track
% ============================================================================
\documentclass{article}

% NeurIPS 2025 style (non-anonymous for D&B track)
\usepackage[final]{neurips_2025}

\usepackage[utf8]{inputenc}
\usepackage[T1]{fontenc}
\usepackage{hyperref}
\usepackage{url}
\usepackage{booktabs}
\usepackage{amsfonts}
\usepackage{amsmath}
\usepackage{amssymb}
\usepackage{nicefrac}
\usepackage{microtype}
\usepackage{graphicx}
\usepackage{xcolor}
\usepackage{algorithm}
\usepackage{algorithmic}
\usepackage{multirow}
\usepackage{subcaption}
\usepackage{enumitem}
\usepackage{tabularx}
\usepackage{makecell}

\definecolor{cblue}{HTML}{4472C4}
\definecolor{cgreen}{HTML}{548235}
\definecolor{corange}{HTML}{ED7D31}

\newcommand{\ours}{\textsc{CrossDisc-Bench}}
\newcommand{\ie}{\textit{i.e.}}
\newcommand{\eg}{\textit{e.g.}}
\newcommand{\cf}{\textit{cf.}}
\newcommand{\etal}{\textit{et al.}}

\title{\ours{}: A Knowledge Graph-Grounded Benchmark for \\
Cross-Disciplinary Scientific Hypothesis Generation}

\author{
  Anonymous Authors \\
  % Affiliation \\
  % \texttt{email@domain.com}
}

\begin{document}
\maketitle

% ============================================================================
\begin{abstract}
% ============================================================================
Cross-disciplinary research---the integration of concepts and methods across distinct scientific fields---has been a driving force behind many major scientific breakthroughs.
However, there exists no systematic benchmark for evaluating the ability of large language models (LLMs) to generate scientifically grounded hypotheses that bridge disciplinary boundaries.
We introduce \ours{}, the first comprehensive benchmark for \emph{cross-disciplinary scientific hypothesis generation} grounded in structured knowledge graphs.
Built from 1{,}000 Nature Communications papers spanning 50 disciplines and 2{,}351 taxonomy terms, \ours{} features:
(1) a multi-stage knowledge extraction pipeline that constructs per-paper knowledge graphs with typed concepts, cross-disciplinary relations, and three-level hypothesis reasoning paths (L1 macro / L2 meso / L3 micro);
(2) an evidence-grounded ground truth (GT) construction method based on constrained term extraction, relation classification, and graph traversal, avoiding reliance on expensive human annotation;
(3) a progressive prompt evaluation framework (P1--P5) that systematically measures how increasing levels of structured knowledge affect hypothesis quality; and
(4) a multi-dimensional evaluation suite combining LLM-as-Judge scoring, text similarity metrics, and novel knowledge graph structural metrics including cross-disciplinary bridging, chain coherence, and Rao--Stirling diversity.
Experiments reveal that knowledge graph-grounded generation (P5) produces 6$\times$ more hypotheses with perfect structural coherence, while free-text methods (P1--P4) achieve higher novelty but lower feasibility.
Our benchmark, code, and data are publicly available to facilitate research on knowledge-intensive scientific discovery.
\end{abstract}

% ============================================================================
\section{Introduction}
\label{sec:intro}
% ============================================================================

Cross-disciplinary research has produced some of the most impactful scientific advances of the past century, from bioinformatics~\cite{lander2001initial} to materials informatics~\cite{agrawal2016perspective}.
The ability to connect concepts across distinct scientific domains---transferring a method from physics to biology, or mapping a mathematical framework onto an engineering problem---lies at the heart of scientific creativity and innovation.
Yet, the cognitive demands of maintaining expertise across multiple fields create substantial barriers to cross-disciplinary discovery~\cite{foster2015tradition}.

The rapid development of large language models (LLMs) has opened new possibilities for automated scientific hypothesis generation~\cite{wang2023scimon,yang2024leandojo}.
Recent benchmarks such as DiscoveryBench~\cite{majumder2024discoverybench} and HypoBench~\cite{liu2025hypobench} have formalized the evaluation of data-driven and classification-based hypothesis generation, respectively.
However, these benchmarks share two critical limitations: (1) they operate primarily within single disciplinary domains and do not assess the ability to bridge distinct fields, and (2) they evaluate hypotheses as free-text outputs without considering the underlying \emph{reasoning structure} that connects concepts across disciplines.

We argue that evaluating cross-disciplinary hypothesis generation requires a fundamentally different approach: one that models the \emph{knowledge graph structure} underlying scientific reasoning.
A high-quality cross-disciplinary hypothesis is not merely a plausible sentence; it is a \emph{structured reasoning path} that starts from a concept in one discipline, traverses typed relations, and arrives at a testable claim in another discipline.

To this end, we introduce \ours{}, the first benchmark that jointly evaluates the quality and structure of cross-disciplinary scientific hypotheses.
Our contributions are as follows:

\begin{enumerate}[leftmargin=*,topsep=2pt,itemsep=1pt]
\item \textbf{CrossDisc-Bench}: A large-scale benchmark constructed from 1{,}000 Nature Communications papers, featuring per-paper knowledge graphs with an average of 28.8 concept nodes, 27.1 typed edges, and three-level hypothesis reasoning paths.
\item \textbf{Evidence-grounded GT construction}: A three-stage pipeline (constrained term extraction $\to$ evidence-based relation classification $\to$ cross-disciplinary graph traversal) that produces structured ground truth without expensive human annotation, yielding an average of 32.0 terms and 21.3 relations per paper.
\item \textbf{Progressive prompt evaluation (P1--P5)}: A systematic framework that incrementally provides structured knowledge (from bare research questions to full knowledge graph context) to isolate the contribution of each knowledge component to hypothesis quality.
\item \textbf{Multi-dimensional evaluation metrics}: A comprehensive suite including LLM-as-Judge scoring across 5 quality dimensions, text similarity metrics, and novel KG structural metrics (bridging score, chain coherence, Rao--Stirling diversity, atypical combination index).
\end{enumerate}

% ============================================================================
\section{Related Work}
\label{sec:related}
% ============================================================================

\paragraph{Scientific Hypothesis Generation.}
Automated hypothesis generation has been studied from multiple perspectives.
DiscoveryBench~\cite{majumder2024discoverybench} formalizes data-driven discovery as $h = \psi(c, v, r)$ over contexts, variables, and relationships, providing 264 real-world tasks across six domains---but focuses on statistical hypothesis verification rather than cross-disciplinary reasoning.
HypoBench~\cite{liu2025hypobench} evaluates LLM-generated hypotheses for classification tasks using Hypothesis Discovery Rate (HDR), with 194 datasets across 12 domains.
SciHyp~\cite{dessi2024scihyp} provides RDF descriptions of 689 hypotheses from computer science papers.
SciMON~\cite{wang2023scimon} generates novel scientific ideas using retrieval-augmented LLMs.
Unlike these works, \ours{} specifically targets \emph{cross-disciplinary} hypothesis generation and evaluates both the \emph{content quality} and the \emph{structural reasoning} behind each hypothesis.

\paragraph{Knowledge Graphs for Scientific Discovery.}
Knowledge graphs (KGs) have been increasingly used to organize and reason over scientific knowledge.
KnowGPT~\cite{luo2024knowgpt} grounds LLM responses in KG facts via graph retrieval-augmented generation.
KG-CoI~\cite{li2024chain} uses knowledge graphs for chain-of-ideas reasoning in hypothesis generation.
In the biomedical domain, systems like BioGPT~\cite{luo2022biogpt} and DrugChat~\cite{liang2023drugchat} leverage domain-specific KGs for drug discovery.
Our work differs by constructing paper-specific knowledge graphs that capture the \emph{cross-disciplinary structure} of individual research papers, rather than relying on pre-existing large-scale KGs.

\paragraph{Cross-Disciplinary Knowledge Integration.}
The challenge of crossing disciplinary boundaries has been studied in scientometrics~\cite{stirling2007general,uzzi2013atypical} and science of science~\cite{foster2015tradition}.
Rao--Stirling diversity~\cite{stirling2007general} and atypical combination indices~\cite{uzzi2013atypical} quantify interdisciplinarity at the journal/citation level.
We adopt and extend these metrics to the knowledge graph level, measuring cross-disciplinary bridging at the concept-relation granularity.

\paragraph{LLM Evaluation Benchmarks.}
The NeurIPS Datasets \& Benchmarks track has seen growing interest in LLM evaluation, from general capabilities~\cite{liang2022holistic} to domain-specific tasks~\cite{sun2024scieval}.
Our work complements this trend by providing the first benchmark specifically designed for evaluating structured, knowledge-grounded scientific reasoning across disciplinary boundaries.

% ============================================================================
\section{Task Formulation}
\label{sec:task}
% ============================================================================

We formally define the task of \emph{cross-disciplinary hypothesis generation} and the key structures involved.

\begin{definition}[Cross-Disciplinary Knowledge Graph]
\label{def:kg}
Given a scientific paper $p$ with title, abstract, and disciplinary classification, a \emph{cross-disciplinary knowledge graph} is a directed attributed graph $\mathcal{G}_p = (\mathcal{V}, \mathcal{E}, \phi_v, \phi_e)$ where:
\begin{itemize}[leftmargin=*,topsep=1pt,itemsep=0pt]
\item $\mathcal{V}$ is a set of concept nodes, each with attributes: term, normalized form, discipline label, evidence text, and confidence score;
\item $\mathcal{E} \subseteq \mathcal{V} \times \mathcal{R} \times \mathcal{V}$ is a set of typed directed edges;
\item $\mathcal{R} = \{$\texttt{method\_applied\_to}, \texttt{improves\_metric}, \texttt{constrains}, \texttt{maps\_to}, \texttt{corresponds\_to}, \texttt{depends\_on}, \texttt{driven\_by}, \texttt{extends}, \texttt{generalizes}, \texttt{inferred\_from}, \texttt{assumes}, \texttt{other}$\}$ is the relation type vocabulary;
\item $\phi_v: \mathcal{V} \to \mathcal{D}$ maps each node to a discipline in taxonomy $\mathcal{T}$;
\item $\phi_e: \mathcal{E} \to [0,1]$ assigns a confidence score to each edge.
\end{itemize}
An edge $(u, r, v) \in \mathcal{E}$ is \emph{cross-disciplinary} if $\phi_v(u) \neq \phi_v(v)$ at the top-level (L1) of taxonomy $\mathcal{T}$.
\end{definition}

\begin{definition}[Hypothesis Reasoning Path]
\label{def:path}
A \emph{hypothesis reasoning path} at granularity level $\ell \in \{L1, L2, L3\}$ is an ordered sequence of three reasoning steps:
\[
\mathbf{h}^\ell = \langle s_1, s_2, s_3 \rangle, \quad s_i = (\text{head}_i, \text{relation}_i, \text{tail}_i, \text{claim}_i)
\]
subject to the \emph{chain consistency} constraint: $\text{tail}_i = \text{head}_{i+1}$ for $i \in \{1, 2\}$.
The final claim $\text{claim}_3$ constitutes the testable hypothesis.
The three granularity levels correspond to: $L1$ (macro-level disciplinary bridging), $L2$ (meso-level mechanism transfer), and $L3$ (micro-level specific predictions).
\end{definition}

\begin{definition}[Cross-Disciplinary Hypothesis Generation Task]
\label{def:task}
Given a paper $p = (\text{title}, \text{abstract}, d_{\text{primary}}, \{d_{\text{secondary}}\})$, the task is to generate a set of hypothesis reasoning paths $\mathcal{H}_p = \{\mathbf{h}_1^{L1}, \ldots, \mathbf{h}_k^{L3}\}$ such that:
\begin{enumerate}[leftmargin=*,topsep=1pt,itemsep=0pt]
\item Each path $\mathbf{h}$ satisfies chain consistency (Def.~\ref{def:path});
\item At least one edge in each path is cross-disciplinary (Def.~\ref{def:kg});
\item The final claim is scientifically grounded, novel, and testable.
\end{enumerate}
\end{definition}

% ============================================================================
\section{\ours{} Construction}
\label{sec:construction}
% ============================================================================

The construction of \ours{} consists of four stages, illustrated in Figure~\ref{fig:pipeline}: data collection and discipline classification (\S\ref{sec:data}), multi-stage knowledge extraction (\S\ref{sec:extraction}), evidence-grounded GT construction (\S\ref{sec:gt}), and dataset statistics (\S\ref{sec:stats}).

\begin{figure}[t]
\centering
\fbox{\parbox{0.95\linewidth}{\centering\small
\textbf{Pipeline Overview} \\[4pt]
1000 Nature Comm. Papers $\xrightarrow{\text{Classification}}$ Cross-Disc. Filter $\xrightarrow{\text{KG Extraction}}$ Per-Paper KGs $\xrightarrow{\text{GT Builder}}$ Benchmark Dataset \\[2pt]
{\scriptsize (OpenAlex $\to$ 3-Level Taxonomy $\to$ Concepts+Relations+Queries+Hypotheses $\to$ Terms+Relations+Paths)}
}}
\caption{Overview of the \ours{} construction pipeline. Each paper undergoes hierarchical discipline classification, multi-stage knowledge extraction to build a structured knowledge graph, and evidence-grounded GT construction via constrained term extraction, relation classification, and graph traversal.}
\label{fig:pipeline}
\end{figure}

% ---------- 4.1 ----------
\subsection{Data Collection and Discipline Classification}
\label{sec:data}

\paragraph{Source data.}
We source papers from the OpenAlex academic database~\cite{priem2022openalex}, selecting English-language articles published in Nature Communications---a high-impact, multidisciplinary journal spanning natural sciences, engineering, and medicine.
From 277{,}720 records in the OpenAlex 2025 snapshot, we filter for Nature Communications papers with abstracts of at least 200 characters, yielding 12{,}082 candidate papers.
We select the top 1{,}000 by citation count (range: XX--XX) and shuffle with a fixed random seed for reproducibility.

\paragraph{Discipline taxonomy.}
We adopt a hierarchical discipline taxonomy derived from the Mathematics Subject Classification (MSC) and extended to cover all natural sciences and engineering fields.
The taxonomy $\mathcal{T}$ contains 50 top-level (L1) disciplines, organized into three levels with 2{,}351 leaf terms.
Examples of L1 disciplines include Chemistry, Materials Science, Biology, Physics, Computer Science, and Clinical Medicine.

\paragraph{Hierarchical classification.}
Each paper is classified into the taxonomy using an LLM-based hierarchical classifier with three rounds of prompting:
(1) L1 classification into top-level disciplines,
(2) L2 refinement within selected L1 branches, and
(3) L3 leaf-level assignment.
A paper is labeled \emph{cross-disciplinary} if it is classified into $\geq$2 distinct L1 disciplines.
The primary discipline is defined as the first L1 classification; all others become secondary disciplines.

In our pilot study of 100 papers, 100\% were classified as cross-disciplinary (expected for Nature Communications).
The most frequent cross-disciplinary pair is Chemistry $\leftrightarrow$ Materials Science (60 papers), followed by Materials Science $\leftrightarrow$ Physics (10) and Clinical Medicine $\leftrightarrow$ Biology (10).

% ---------- 4.2 ----------
\subsection{Multi-Stage Knowledge Extraction}
\label{sec:extraction}

For each cross-disciplinary paper, we run a four-stage knowledge extraction pipeline using LLM calls with structured output parsing (Figure~\ref{fig:pipeline}).

\paragraph{Stage 1: Concept extraction.}
We extract domain-specific concepts organized by discipline.
Each concept is represented as a \texttt{ConceptEntry} with fields: term, normalized form, standard label (grounded to $\mathcal{T}$), evidence text ($\leq$40 characters), source, and confidence $\in [0,1]$.
Concepts are partitioned into primary discipline concepts and secondary discipline concepts.
On average, we extract $15.6 \pm 2.9$ concepts per paper ($13.4$ primary, $2.2$ secondary).

\paragraph{Stage 2: Cross-disciplinary relation extraction.}
We identify typed relations between concepts using the 12-type relation vocabulary $\mathcal{R}$ (Def.~\ref{def:kg}).
Each relation is a \texttt{RelationEntry} with head, tail, relation type, direction, quantitative information (if applicable), assumptions, evidence, and confidence.
We extract an average of $8.1 \pm 1.3$ relations per paper, totaling 765 relations across 94 papers.
Table~\ref{tab:relation_dist} shows the distribution of relation types.

\begin{table}[t]
\centering
\caption{Distribution of extracted cross-disciplinary relation types (94 papers).}
\label{tab:relation_dist}
\small
\begin{tabular}{lrr}
\toprule
\textbf{Relation Type} & \textbf{Count} & \textbf{\%} \\
\midrule
\texttt{method\_applied\_to} & 225 & 29.4\% \\
\texttt{improves\_metric} & 166 & 21.7\% \\
\texttt{constrains} & 104 & 13.6\% \\
\texttt{maps\_to} & 62 & 8.1\% \\
\texttt{corresponds\_to} & 48 & 6.3\% \\
\texttt{depends\_on} & 41 & 5.4\% \\
\texttt{driven\_by} & 38 & 5.0\% \\
\texttt{extends} & 32 & 4.2\% \\
\texttt{other} & 49 & 6.4\% \\
\bottomrule
\end{tabular}
\end{table}

\paragraph{Stage 3: Hierarchical query generation.}
We generate research queries at three granularity levels aligned with the hypothesis path structure:
\begin{itemize}[leftmargin=*,topsep=1pt,itemsep=0pt]
\item \textbf{L1 (macro):} A single broad question about the paper's core cross-disciplinary contribution.
\item \textbf{L2 (meso):} 2--3 questions about specific mechanism transfers or methodological bridges.
\item \textbf{L3 (micro):} 2--3 questions about concrete, testable predictions.
\end{itemize}

\paragraph{Stage 4: Three-level hypothesis path generation.}
For each query level, we generate structured hypothesis reasoning paths (Def.~\ref{def:path}).
Each path consists of exactly three reasoning steps with enforced chain consistency ($\text{tail}_i = \text{head}_{i+1}$).
The generation is parallelized across L1, L2, and L3 using concurrent API calls.
On average, we obtain $2.5$ L1 paths, $2.0$ L2 paths, and $2.0$ L3 paths per paper, each with 3 steps.

\paragraph{Knowledge graph construction.}
After extraction, we assemble a per-paper knowledge graph $\mathcal{G}_p$ by:
(1) creating nodes from extracted concepts (with discipline labels from taxonomy grounding),
(2) creating edges from both structural relations and hypothesis path steps,
(3) computing graph-level quality metrics.
The resulting graphs have an average of $28.8 \pm 6.0$ nodes and $27.1 \pm 2.9$ edges with a density of $0.076 \pm 0.030$.

% ---------- 4.3 ----------
\subsection{Evidence-Grounded GT Construction}
\label{sec:gt}

A key contribution of \ours{} is an \emph{evidence-grounded} approach to GT construction that avoids the scalability bottleneck of manual annotation.
Our GT builder operates in three stages:

\begin{algorithm}[t]
\caption{Evidence-Grounded GT Construction}
\label{alg:gt}
\begin{algorithmic}[1]
\REQUIRE Paper $p$ with title, abstract, extracted concepts $C_p$, taxonomy $\mathcal{T}$
\ENSURE Ground truth $\text{GT}_p = (\text{terms}, \text{relations}, \text{paths})$
\STATE \textbf{Stage 1: Constrained Term Extraction}
\STATE $\text{terms} \leftarrow \text{GroundTerms}(C_p, \mathcal{T})$ \COMMENT{Ground concepts against taxonomy}
\FOR{each term $t \in \text{terms}$}
    \STATE Assign discipline $\phi(t)$, evidence, confidence
\ENDFOR
\STATE \textbf{Stage 2: Evidence-Based Relation Construction}
\STATE Split abstract into sentences $S = \{s_1, \ldots, s_n\}$
\FOR{each sentence $s_j$ containing terms $t_a, t_b \in \text{terms}$}
    \STATE $r_{ab} \leftarrow \text{ClassifyRelation}(t_a, t_b, s_j)$
    \STATE Add $(t_a, r_{ab}, t_b, s_j)$ to relations with evidence
\ENDFOR
\STATE \textbf{Stage 3: Cross-Disciplinary Path Construction}
\STATE Build directed graph $G$ from terms (nodes) and relations (edges)
\FOR{each pair $(u, v)$ where $\phi(u) \neq \phi(v)$}
    \STATE Find all simple paths $u \leadsto v$ with length $\leq 4$
\ENDFOR
\STATE Rank paths by (\# disciplines crossed, evidence confidence)
\RETURN top-$k$ paths as GT hypothesis paths
\end{algorithmic}
\end{algorithm}

\paragraph{Stage 1: Constrained term extraction.}
We leverage the high-quality concepts already extracted by the production pipeline (\S\ref{sec:extraction}) and ground them against the taxonomy $\mathcal{T}$.
Terms are augmented with heuristic extraction of acronyms, scientific phrases, and hyphenated compounds from the abstract.
Each term receives a grounding confidence score based on string similarity to taxonomy entries.

\paragraph{Stage 2: Evidence-based relation construction.}
We identify term co-occurrences within individual sentences of the abstract and classify each co-occurrence pair into the 12-type relation vocabulary using either LLM-based classification or keyword-pattern heuristics.
Each relation is accompanied by the evidence sentence from which it was derived, ensuring traceability.

\paragraph{Stage 3: Cross-disciplinary path construction.}
We construct a directed graph from terms and relations, then enumerate all simple paths between nodes of different L1 disciplines using NetworkX~\cite{hagberg2008exploring}.
Paths are ranked by the number of disciplines crossed and the cumulative evidence confidence, and the top-$k$ ($k=20$) paths are retained as GT hypothesis paths.

This evidence-grounded approach yields GT entries with an average of $32.0 \pm 6.7$ terms and $21.3 \pm 14.1$ relations per paper (see Table~\ref{tab:gt_stats}).

% ---------- 4.4 ----------
\subsection{Dataset Statistics}
\label{sec:stats}

\begin{table}[t]
\centering
\caption{Summary statistics of \ours{} (pilot: 94 papers; full: 1{,}000 papers).}
\label{tab:dataset_stats}
\small
\begin{tabular}{lcc}
\toprule
\textbf{Statistic} & \textbf{Pilot (94)} & \textbf{Full (1{,}000)} \\
\midrule
\multicolumn{3}{l}{\textit{Knowledge Graph}} \\
Avg. concept nodes / paper & $28.8 \pm 6.0$ & --- \\
Avg. typed edges / paper & $27.1 \pm 2.9$ & --- \\
Total concept nodes & 2{,}711 & --- \\
Total typed edges & 2{,}549 & --- \\
Graph density & $0.076 \pm 0.030$ & --- \\
Unique relation types & 12 & 12 \\
\midrule
\multicolumn{3}{l}{\textit{Hypothesis Paths}} \\
Avg. L1 paths / paper & $2.5 \pm 0.5$ & --- \\
Avg. L2 paths / paper & $2.0 \pm 0.1$ & --- \\
Avg. L3 paths / paper & $2.0 \pm 0.4$ & --- \\
Steps per path & 3 (fixed) & 3 (fixed) \\
Chain consistency rate & 100\% & --- \\
\midrule
\multicolumn{3}{l}{\textit{Benchmark GT}} \\
GT entries & 91 & --- \\
Avg. terms / entry & $32.0 \pm 6.7$ & --- \\
Avg. relations / entry & $21.3 \pm 14.1$ & --- \\
Avg. cross-disc. paths / entry & $11.7$ & --- \\
\midrule
\multicolumn{3}{l}{\textit{Discipline Coverage}} \\
L1 disciplines represented & 50 & 50 \\
Leaf taxonomy terms & 2{,}351 & 2{,}351 \\
Top cross-disc. pair & Chem $\leftrightarrow$ MatSci & --- \\
\bottomrule
\end{tabular}
\end{table}

Table~\ref{tab:dataset_stats} summarizes the key statistics.
The pilot study (94 papers) demonstrates the pipeline's effectiveness, with a 94\% extraction success rate (94/100 papers).
The top-5 primary disciplines are Chemistry (53\%), Materials Science (14\%), Biology (12\%), Physics (6\%), and Computer Science (6\%), reflecting Nature Communications' broad coverage.

\begin{table}[t]
\centering
\caption{Knowledge graph quality metrics across 94 papers (mean $\pm$ std).}
\label{tab:gt_stats}
\small
\begin{tabular}{lcc}
\toprule
\textbf{Metric} & \textbf{Mean $\pm$ Std} & \textbf{Range} \\
\midrule
\multicolumn{3}{l}{\textit{Core Quality}} \\
Path Consistency & $0.203 \pm 0.121$ & [0.00, 0.52] \\
Concept Coverage & $0.481 \pm 0.133$ & [0.18, 0.82] \\
Bridging Score & $0.390 \pm 0.180$ & [0.05, 0.85] \\
Chain Coherence & $0.580 \pm 0.052$ & [0.44, 0.72] \\
\midrule
\multicolumn{3}{l}{\textit{Diversity \& Novelty}} \\
Rao--Stirling Diversity & $0.141 \pm 0.089$ & [0.00, 0.45] \\
Embedding Bridging & $0.962 \pm 0.031$ & [0.87, 1.00] \\
Atypical Combination & $0.312 \pm 0.156$ & [0.00, 0.71] \\
\midrule
\multicolumn{3}{l}{\textit{Topology}} \\
Graph Density & $0.076 \pm 0.030$ & [0.02, 0.16] \\
Clustering Coefficient & $0.184 \pm 0.098$ & [0.00, 0.45] \\
Largest CC Ratio & $0.827 \pm 0.112$ & [0.48, 1.00] \\
Avg. Betweenness & $0.042 \pm 0.021$ & [0.01, 0.12] \\
\bottomrule
\end{tabular}
\end{table}

% ============================================================================
\section{Evaluation Framework}
\label{sec:eval}
% ============================================================================

% ---------- 5.1 ----------
\subsection{Progressive Prompt Levels (P1--P5)}
\label{sec:prompts}

To systematically measure the effect of structured knowledge on hypothesis quality, we design five prompt levels with progressively increasing information:

\begin{itemize}[leftmargin=*,topsep=2pt,itemsep=2pt]
\item \textbf{P1 (Bare Question):} Only the L1 research query is provided. No paper metadata, no extracted knowledge. The LLM must generate hypotheses purely from parametric knowledge.
\item \textbf{P2 (+ Abstract \& Disciplines):} Adds the paper's abstract and discipline role (primary/secondary classification). Provides domain context without structured knowledge.
\item \textbf{P3 (+ Concepts):} Adds the list of extracted domain-specific concepts organized by discipline. Enables concept-aware reasoning.
\item \textbf{P4 (+ Relations):} Adds known cross-disciplinary relations between concepts. Provides explicit structural information about how concepts connect across fields.
\item \textbf{P5 (Full KG):} Uses the complete knowledge graph including concepts, relations, queries at all three levels, and the structured extraction output. Generates hypothesis paths with strict 3-step chain consistency constraints.
\end{itemize}

This progressive design isolates the marginal contribution of each knowledge component: P1$\to$P2 measures the effect of domain context, P2$\to$P3 measures concept awareness, P3$\to$P4 measures relational structure, and P4$\to$P5 measures full KG grounding.

% ---------- 5.2 ----------
\subsection{Multi-Dimensional Evaluation Metrics}
\label{sec:metrics}

We evaluate generated hypotheses along three complementary dimensions.

\paragraph{Dimension 1: LLM-as-Judge quality scoring.}
Following recent evaluation practices~\cite{zheng2024judging}, we use an LLM judge to rate each hypothesis on five dimensions (1--10 scale):
\begin{itemize}[leftmargin=*,topsep=1pt,itemsep=0pt]
\item \textbf{Novelty}: Does the hypothesis go beyond the paper's original conclusions?
\item \textbf{Specificity}: Is the hypothesis concrete, testable, and operationally defined?
\item \textbf{Feasibility}: Can the hypothesis be verified with existing experimental technology?
\item \textbf{Relevance}: Is the hypothesis aligned with the paper's research direction?
\item \textbf{Cross-disciplinarity}: Does the hypothesis meaningfully integrate multiple disciplines?
\end{itemize}

\paragraph{Dimension 2: Text similarity metrics.}
We compute standard text similarity metrics between generated hypotheses and reference summaries:
BERTScore~\cite{zhang2019bertscore} (using DeBERTa-xlarge-MNLI), ROUGE-1/2/L~\cite{lin2004rouge}, and BLEU~\cite{papineni2002bleu}.
Note that low similarity to references is \emph{expected} for novel hypotheses; these metrics primarily measure relevance rather than quality.

\paragraph{Dimension 3: Knowledge graph structural metrics.}
For structured outputs (primarily P5), we compute novel graph-level metrics:

\begin{itemize}[leftmargin=*,topsep=1pt,itemsep=0pt]
\item \textbf{Chain Coherence}: The fraction of consecutive step pairs $(s_i, s_{i+1})$ satisfying $\text{tail}_i = \text{head}_{i+1}$, measuring reasoning path validity:
\begin{equation}
\text{CC}(\mathbf{h}) = \frac{1}{|\mathbf{h}|-1} \sum_{i=1}^{|\mathbf{h}|-1} \mathbb{1}[\text{tail}_i = \text{head}_{i+1}]
\end{equation}
\item \textbf{Cross-Disciplinary Bridging}: The proportion of edges connecting nodes from different L1 disciplines:
\begin{equation}
\text{Bridge}(\mathcal{G}) = \frac{|\{(u,r,v) \in \mathcal{E} : \phi_v(u) \neq \phi_v(v)\}|}{|\mathcal{E}_{\text{valid}}|}
\end{equation}
\item \textbf{Rao--Stirling Diversity}~\cite{stirling2007general}: $\Delta = \sum_{i \neq j} d_{ij} \cdot p_i \cdot p_j$, where $p_i$ is the proportion of concepts in discipline $i$ and $d_{ij}$ is the semantic distance between disciplines.
\item \textbf{Entity Grounding Rate}: The fraction of head/tail entities in hypothesis paths that are grounded to specific scientific terms (not generic phrases).
\item \textbf{Relation Diversity}: The number of unique relation types divided by total paths, measuring the variety of reasoning patterns.
\end{itemize}

% ============================================================================
\section{Experiments}
\label{sec:experiments}
% ============================================================================

% ---------- 6.1 ----------
\subsection{Experimental Setup}
\label{sec:setup}

\paragraph{Models.}
We use Qwen3-235B-A22B as the primary LLM for both knowledge extraction and hypothesis generation.
For evaluation, we use the same model as the LLM-as-Judge.
All API calls use streaming to avoid gateway timeouts, with exponential backoff retry (up to 5 attempts, wait 5--120s) for network errors and rate limiting.

\paragraph{Parameters.}
Knowledge extraction uses temperature 0.2 with seed 42 for reproducibility.
Hypothesis generation uses temperature 0.3 with max\_tokens=8192 to avoid truncation.
Query generation uses temperature 0.2.
The batch extraction pipeline uses 8 parallel workers.

\paragraph{Test protocol.}
From the successfully extracted papers, we reserve the last 10 as the test set and use the remaining as the GT construction set.
For each test paper, we generate hypotheses at all five prompt levels (P1--P5) and evaluate using all three metric dimensions.

% ---------- 6.2 ----------
\subsection{Main Results: P1--P5 Comparison}
\label{sec:main_results}

\begin{table}[t]
\centering
\caption{P1--P5 hypothesis generation comparison (mean across test papers). \textbf{Bold}: best per column. \underline{Underline}: second best. $^\dagger$P5 directly uses extracted structured paths.}
\label{tab:p1p5}
\small
\begin{tabular}{lccccccc}
\toprule
\textbf{Method} & \textbf{Nov.} & \textbf{Spec.} & \textbf{Feas.} & \textbf{Relev.} & \textbf{Cross-D.} & \textbf{\#Hyp} & \textbf{Time (s)} \\
\midrule
P1 (Bare) & \textbf{8.67} & \textbf{9.00} & 8.00 & 9.33 & \textbf{9.00} & 1.0 & 22.8 \\
P2 (+Abs) & 8.00 & \underline{8.67} & \underline{8.33} & \textbf{9.67} & 8.33 & 1.0 & 31.9 \\
P3 (+Con) & 8.00 & 8.33 & 8.00 & 9.33 & \underline{8.67} & 1.0 & 33.7 \\
P4 (+Rel) & --- & --- & --- & --- & --- & 1.0 & 51.9 \\
P5$^\dagger$ (Full KG) & 7.67 & 7.33 & \textbf{8.67} & 9.33 & 8.00 & \textbf{6.33} & $<$0.01 \\
\bottomrule
\end{tabular}
\end{table}

Table~\ref{tab:p1p5} presents the main comparison across prompt levels.
Several key observations emerge:

\paragraph{Novelty-feasibility trade-off.}
P1 (bare question, no context) achieves the highest novelty (8.67) and cross-disciplinarity (9.00), as the LLM generates creative hypotheses unconstrained by paper-specific knowledge.
However, P5 (full KG) achieves the highest feasibility (8.67), as structured knowledge grounds hypotheses in established concepts and relations.
This trade-off between unconstrained creativity and knowledge-grounded precision is a central finding of our benchmark.

\paragraph{Information does not monotonically improve quality.}
Contrary to the intuition that ``more information is always better,'' we observe that quality metrics do not monotonically increase from P1 to P4.
P2 achieves the highest relevance (9.67) by providing domain context, but P3 and P4 show slight decreases in novelty, suggesting that explicit concept lists may constrain the LLM's creative space.

\paragraph{Structured vs.\ free-text generation.}
P5 generates $6.3\times$ more hypotheses than P1--P4 (6.33 vs.\ 1.0) by leveraging the multi-level path structure.
The generated paths achieve perfect chain coherence (1.0) and entity grounding (1.0) by construction.
However, free-text methods produce more diverse and novel individual hypotheses.

% ---------- 6.3 ----------
\subsection{KG Quality Analysis}
\label{sec:kg_quality}

Table~\ref{tab:gt_stats} reports the knowledge graph quality metrics across all 94 papers.
The moderate concept coverage (0.481) indicates that hypothesis paths cover roughly half of the knowledge graph's concept space, leaving room for alternative reasoning paths.
The bridging score of 0.390 confirms that a substantial fraction of edges cross disciplinary boundaries, validating the cross-disciplinary nature of the extracted knowledge.
High embedding bridging (0.962) indicates that hypothesis paths connect semantically distant concepts, a hallmark of genuine cross-disciplinary reasoning.

% ---------- 6.4 ----------
\subsection{Depth Analysis: L1 vs.\ L2 vs.\ L3}
\label{sec:depth}

\begin{table}[t]
\centering
\caption{KG-based evaluation by hypothesis depth level (P5, mean $\pm$ std across 3 test papers).}
\label{tab:depth}
\small
\begin{tabular}{lcccc}
\toprule
\textbf{Metric} & \textbf{L1 (Macro)} & \textbf{L2 (Meso)} & \textbf{L3 (Micro)} \\
\midrule
Bridging Score & $1.000 \pm 0.000$ & $0.891 \pm 0.105$ & $0.756 \pm 0.142$ \\
Chain Coherence & $0.559 \pm 0.060$ & $0.599 \pm 0.048$ & $0.570 \pm 0.047$ \\
Innovation (LLM) & $7.00 \pm 1.30$ & $6.87 \pm 0.51$ & $6.92 \pm 0.31$ \\
Feasibility (LLM) & $7.06 \pm 0.57$ & $7.42 \pm 0.12$ & $7.05 \pm 0.87$ \\
Scientificity (LLM) & $6.92 \pm 1.23$ & $7.92 \pm 0.12$ & $7.45 \pm 0.27$ \\
Testability (LLM) & $7.68 \pm 0.30$ & $7.81 \pm 0.33$ & $7.77 \pm 0.28$ \\
\midrule
Concept Expansion & --- & $0.268 \pm 0.191$ & $0.472 \pm 0.239$ \\
Anchoring & --- & $0.621 \pm 0.039$ & $0.387 \pm 0.169$ \\
Depth Quality & --- & \multicolumn{2}{c}{$0.216 \pm 0.125$} \\
\bottomrule
\end{tabular}
\end{table}

Table~\ref{tab:depth} reveals interesting patterns across hypothesis depth levels:

\begin{itemize}[leftmargin=*,topsep=1pt,itemsep=1pt]
\item \textbf{Bridging decreases with depth}: L1 paths have perfect bridging (1.0) as they explicitly connect primary and secondary disciplines. L2 and L3 paths operate within more specific subdomains, naturally reducing cross-disciplinary edge proportion.
\item \textbf{Scientificity increases at L2}: L2 (meso-level) paths achieve the highest scientificity score (7.92), suggesting that intermediate-granularity hypotheses best balance specificity with scientific rigor.
\item \textbf{Concept expansion increases with depth}: L3 paths expand the concept space by 47.2\% beyond L1, while L2 expands by 26.8\%, indicating that deeper reasoning introduces more specialized vocabulary.
\item \textbf{Anchoring decreases with depth}: L2 paths are well-anchored to L1 concepts (0.621), but L3 paths diverge further (0.387), introducing novel entities not present at higher levels.
\end{itemize}

% ---------- 6.5 ----------
\subsection{Case Study}
\label{sec:case}

We illustrate the benchmark with a paper on nano-island-encapsulated cobalt single-atom catalysts for Fenton-like reactions, classified as cross-disciplinary between Chemistry (primary) and Materials Science / Environmental Science (secondary).

\paragraph{Extracted KG.}
The knowledge graph contains 31 concept nodes across three disciplines and 28 typed edges.
Key cross-disciplinary relations include:
\emph{nano-island encapsulation} $\xrightarrow{\texttt{method\_applied\_to}}$ \emph{Co single atoms} (Chemistry$\to$Materials Science),
\emph{PMS activation} $\xrightarrow{\texttt{corresponds\_to}}$ \emph{sulfate radicals} (Chemistry$\to$Environmental Science).

\paragraph{L1 hypothesis path.}
Step 1: \emph{Fenton-like catalysis} faces \emph{activity-stability trade-off} [Chemistry bottleneck].
Step 2: \emph{Activity-stability trade-off} is addressed by \emph{nano-island encapsulation design} [Materials Science solution].
Step 3: \emph{Nano-island encapsulation} enables \emph{precise electronic state control of cobalt single atoms}, improving catalytic efficiency while maintaining stability [testable claim].

This three-step path bridges Chemistry and Materials Science through a mechanism transfer, illustrating the type of cross-disciplinary reasoning our benchmark is designed to evaluate.

% ============================================================================
\section{Limitations and Broader Impact}
\label{sec:limitations}
% ============================================================================

\paragraph{Limitations.}
Our work has several limitations that we acknowledge transparently:
\begin{itemize}[leftmargin=*,topsep=1pt,itemsep=1pt]
\item \textbf{LLM dependency}: Both extraction and evaluation rely on LLM calls, introducing potential biases from the underlying model. We mitigate this by using structured output parsing with Pydantic schema validation and evidence-grounded GT construction.
\item \textbf{Discipline coverage bias}: Nature Communications skews toward natural sciences and engineering. Social sciences, humanities, and formal sciences are underrepresented.
\item \textbf{English only}: Our current pipeline processes English-language papers only, though the taxonomy includes bilingual (English/Chinese) term mappings.
\item \textbf{GT quality bounds}: While evidence-grounded GT avoids annotation bottlenecks, it cannot capture all valid hypothesis paths. Multiple valid cross-disciplinary hypotheses may exist for any given paper.
\item \textbf{Single-model baseline}: Our current experiments use a single LLM (Qwen3-235B). Multi-model comparisons (GPT-4o, DeepSeek-R1, Llama-3, etc.) are needed for robust conclusions.
\end{itemize}

\paragraph{Broader impact.}
\ours{} is designed to advance research on automated scientific discovery, which has significant positive potential for accelerating cross-disciplinary breakthroughs.
However, we caution against using generated hypotheses as substitutes for rigorous scientific methodology.
LLM-generated hypotheses should be treated as \emph{starting points} that require expert validation, experimental testing, and peer review before informing research directions.
We do not foresee direct negative societal impacts from this benchmark, as it evaluates scientific reasoning capabilities rather than generating actionable claims in sensitive domains.

% ============================================================================
\section{Conclusion}
\label{sec:conclusion}
% ============================================================================

We have presented \ours{}, the first comprehensive benchmark for cross-disciplinary scientific hypothesis generation grounded in structured knowledge graphs.
Built from 1{,}000 Nature Communications papers across 50 disciplines, \ours{} introduces several novel contributions: an evidence-grounded GT construction pipeline that scales without human annotation, a progressive prompt evaluation framework (P1--P5) that isolates the contribution of structured knowledge, and a multi-dimensional evaluation suite combining LLM-as-Judge scoring with knowledge graph structural metrics.

Our experiments reveal a fundamental \emph{novelty-feasibility trade-off}: unconstrained generation (P1) produces the most creative hypotheses, while knowledge-grounded generation (P5) produces the most structurally coherent and feasible ones.
This finding suggests that future systems should combine the creative strengths of free-text generation with the structural guarantees of knowledge graph grounding.

We release \ours{} as an open resource to the community, including the complete dataset, extraction pipeline code, and evaluation framework, to facilitate further research on knowledge-intensive cross-disciplinary scientific discovery.

% ============================================================================
% References
% ============================================================================
\bibliographystyle{plain}
% \bibliography{references}

\begin{thebibliography}{20}

\bibitem{agrawal2016perspective}
A.~Agrawal and A.~Choudhary.
\newblock Perspective: Materials informatics and big data: Realization of the ``fourth paradigm'' of science in materials science.
\newblock {\em APL Materials}, 4(5):053208, 2016.

\bibitem{dessi2024scihyp}
D.~Dess\`{\i} \etal{}
\newblock SciHyp: A fine-grained dataset describing hypotheses and their components from scientific articles.
\newblock In {\em ISWC}, 2024.

\bibitem{foster2015tradition}
J.~G. Foster, A.~Rzhetsky, and J.~A. Evans.
\newblock Tradition and innovation in scientists' research strategies.
\newblock {\em American Sociological Review}, 80(5):875--908, 2015.

\bibitem{hagberg2008exploring}
A.~Hagberg, P.~Swart, and D.~Chult.
\newblock Exploring network structure, dynamics, and function using NetworkX.
\newblock Technical report, Los Alamos National Lab, 2008.

\bibitem{lander2001initial}
E.~S. Lander \etal{}
\newblock Initial sequencing and analysis of the human genome.
\newblock {\em Nature}, 409(6822):860--921, 2001.

\bibitem{li2024chain}
Z.~Li \etal{}
\newblock Chain of ideas: Revolutionizing research via novel idea development with LLM agents.
\newblock {\em arXiv preprint arXiv:2410.13185}, 2024.

\bibitem{liang2022holistic}
P.~Liang \etal{}
\newblock Holistic evaluation of language models.
\newblock {\em Transactions on Machine Learning Research}, 2022.

\bibitem{liang2023drugchat}
Y.~Liang \etal{}
\newblock DrugChat: Towards enabling ChatGPT-like capabilities on drug molecule graphs.
\newblock {\em arXiv preprint arXiv:2309.16712}, 2023.

\bibitem{lin2004rouge}
C.-Y. Lin.
\newblock ROUGE: A package for automatic evaluation of summaries.
\newblock In {\em Text Summarization Branches Out}, 2004.

\bibitem{liu2025hypobench}
H.~Liu \etal{}
\newblock HypoBench: Towards systematic and principled benchmarking for hypothesis generation.
\newblock {\em arXiv preprint arXiv:2504.11524}, 2025.

\bibitem{luo2024knowgpt}
L.~Luo \etal{}
\newblock KnowGPT: Knowledge graph based prompting for large language models.
\newblock In {\em NeurIPS}, 2024.

\bibitem{luo2022biogpt}
R.~Luo \etal{}
\newblock BioGPT: Generative pre-trained transformer for biomedical text generation and mining.
\newblock {\em Briefings in Bioinformatics}, 23(6):bbac409, 2022.

\bibitem{majumder2024discoverybench}
B.~P. Majumder \etal{}
\newblock DiscoveryBench: Towards data-driven discovery with large language models.
\newblock In {\em NeurIPS Datasets and Benchmarks Track}, 2024.

\bibitem{papineni2002bleu}
K.~Papineni \etal{}
\newblock BLEU: A method for automatic evaluation of machine translation.
\newblock In {\em ACL}, 2002.

\bibitem{priem2022openalex}
J.~Priem, H.~Piwowar, and R.~Orr.
\newblock OpenAlex: A fully-open index of scholarly works, authors, venues, institutions, and concepts.
\newblock {\em arXiv preprint arXiv:2205.01833}, 2022.

\bibitem{stirling2007general}
A.~Stirling.
\newblock A general framework for analysing diversity in science, technology and society.
\newblock {\em Journal of the Royal Society Interface}, 4(15):707--719, 2007.

\bibitem{sun2024scieval}
L.~Sun \etal{}
\newblock SciEval: A multi-level large language model evaluation benchmark for scientific research.
\newblock In {\em AAAI}, 2024.

\bibitem{uzzi2013atypical}
B.~Uzzi \etal{}
\newblock Atypical combinations and scientific impact.
\newblock {\em Science}, 342(6157):468--472, 2013.

\bibitem{wang2023scimon}
H.~Wang \etal{}
\newblock SciMON: Scientific inspiration machines optimized for novelty.
\newblock In {\em ACL}, 2023.

\bibitem{yang2024leandojo}
K.~Yang \etal{}
\newblock LeanDojo: Theorem proving with retrieval-augmented language models.
\newblock In {\em NeurIPS}, 2024.

\bibitem{zhang2019bertscore}
T.~Zhang \etal{}
\newblock BERTScore: Evaluating text generation with BERT.
\newblock In {\em ICLR}, 2020.

\bibitem{zheng2024judging}
L.~Zheng \etal{}
\newblock Judging LLM-as-a-judge with MT-Bench and Chatbot Arena.
\newblock In {\em NeurIPS}, 2024.

\end{thebibliography}

% ============================================================================
% Appendix
% ============================================================================
\newpage
\appendix

\section{Prompt Templates}
\label{app:prompts}

We provide the core structure of our five prompt levels. Full prompts are available in the supplementary code.

\paragraph{P1 System Prompt (excerpt).}
\begin{quote}
\small\itshape
You are a cross-disciplinary research hypothesis generator. Given a research question, generate a novel, testable hypothesis that bridges at least two scientific disciplines. The hypothesis should be specific enough to be experimentally verifiable.
\end{quote}

\paragraph{P2 additions.}
Adds: paper abstract, primary discipline name, secondary discipline names, and their roles in the research.

\paragraph{P3 additions.}
Adds: extracted concept list organized by discipline, with each concept's evidence text and confidence.

\paragraph{P4 additions.}
Adds: cross-disciplinary relation list with head/tail concepts, relation types, and evidence. Explicit instruction to reason along the provided relational structure.

\paragraph{P5 additions.}
Uses the full structured extraction output to generate three-step hypothesis reasoning paths at L1/L2/L3 granularity, with strict chain consistency constraints ($\text{tail}_i = \text{head}_{i+1}$).

\section{Datasheet for \ours{}}
\label{app:datasheet}

Following the Datasheets for Datasets framework~\cite{gebru2021datasheets}, we provide key documentation:

\paragraph{Motivation.}
\ours{} was created to fill the gap in benchmarks for evaluating cross-disciplinary scientific hypothesis generation. It is intended for researchers working on LLM-based scientific discovery, knowledge graph reasoning, and evaluation methodology.

\paragraph{Composition.}
The dataset contains structured knowledge graphs and hypothesis paths for 1{,}000 Nature Communications papers. Each instance includes: paper metadata (title, abstract, DOI, disciplines), extracted knowledge graph (concepts, relations), three-level hypothesis paths (L1/L2/L3), and evidence-grounded GT (terms, relations, cross-disciplinary paths).

\paragraph{Collection process.}
Paper metadata was collected from OpenAlex (CC0 license). Knowledge extraction and GT construction were performed using LLM-based pipelines with structured output validation. No human subjects were involved.

\paragraph{Uses.}
Intended uses include: benchmarking LLMs on cross-disciplinary reasoning, developing knowledge graph-grounded generation methods, and studying the relationship between structured knowledge and hypothesis quality. The dataset should \emph{not} be used to make clinical or policy decisions based on generated hypotheses.

\paragraph{Distribution.}
The dataset will be distributed via Hugging Face Datasets under the CC-BY-4.0 license. Code is available on GitHub under the MIT license.

\paragraph{Maintenance.}
The dataset will be maintained by the authors. We plan to update the benchmark with additional papers and models as they become available.

\section{Additional Experimental Results}
\label{app:results}

\begin{table}[h]
\centering
\caption{Discipline distribution in the benchmark dataset (top 10 primary disciplines).}
\label{tab:disc_dist}
\small
\begin{tabular}{lrr}
\toprule
\textbf{Primary Discipline} & \textbf{Count} & \textbf{\%} \\
\midrule
Chemistry & 53 & 53.0\% \\
Materials Science & 14 & 14.0\% \\
Biology & 12 & 12.0\% \\
Physics & 6 & 6.0\% \\
Computer Science & 6 & 6.0\% \\
Clinical Medicine & 4 & 4.0\% \\
Environmental Science & 2 & 2.0\% \\
Basic Medicine & 1 & 1.0\% \\
Electronics \& Communications & 1 & 1.0\% \\
Energy Science & 1 & 1.0\% \\
\bottomrule
\end{tabular}
\end{table}

\begin{table}[h]
\centering
\caption{Cross-disciplinary pair frequency (top 10 pairs).}
\label{tab:cross_pairs}
\small
\begin{tabular}{lr}
\toprule
\textbf{Discipline Pair} & \textbf{Count} \\
\midrule
Chemistry $\leftrightarrow$ Materials Science & 60 \\
Materials Science $\leftrightarrow$ Physics & 10 \\
Clinical Medicine $\leftrightarrow$ Biology & 10 \\
Chemistry $\leftrightarrow$ Physics & 8 \\
Biology $\leftrightarrow$ Computer Science & 8 \\
Chemistry $\leftrightarrow$ Energy Science & 7 \\
Chemistry $\leftrightarrow$ Environmental Science & 5 \\
Basic Medicine $\leftrightarrow$ Biology & 5 \\
Electronics $\leftrightarrow$ Physics & 4 \\
Mechanical Eng. $\leftrightarrow$ Materials Science & 3 \\
\bottomrule
\end{tabular}
\end{table}

\end{document}
